\section{Introduction}

Large language models have moved quickly from answering medical examination questions to acting
inside clinical workflows. Systems now select and call validated clinical tools
\cite{gao2025txagent,liu2025riskagent}, conduct multi-turn diagnostic conversations
\cite{palepu2025disease,saab2025multimodal}, and in at least one case operate prospectively at the
bedside \cite{shashikumar2025sepsis}. The building blocks of such systems --- ReAct-style reasoning
\cite{yao2023react}, retrieval-augmented generation \cite{lewis2020rag}, multi-agent coordination
\cite{wu2024autogen}, and structured agent memory \cite{wu2025memory} --- are individually mature.

What has not matured is the evidence that these systems can be trusted with a specific patient
over time. Three shortfalls recur across the strongest published work. First, evaluation remains
distant from deployment: medical agents are still graded on curated examinations or synthetic
records \cite{tang2024medagents,li2024agenthospital,jiang2025medagentbench}, not on the noisy,
longitudinal, internally revised record of a real intensive-care admission \cite{johnson2023mimic}.
Second, retrieval is grounded in guidelines and literature rather than in the patient: even
strong clinical RAG deployments treat external knowledge as the corpus and the patient as the
query \cite{zhao2025medrag,ke2025ragfitness}, leaving the patient's own timeline --- prior labs, prior
admissions, the trend of a vital sign --- outside the evidence base. Third, verification and
auditability are asserted rather than measured: hallucination and uncertainty are now well
characterized \cite{kim2025hallucinations,atf2025uncertainty}, and regulators increasingly demand
inspectable safeguards for agentic clinical software \cite{tan2026undcs,weissman2025unregulated},
yet published systems rarely check a recommendation against the retrieved patient evidence, and
none that we are aware of measures the faithfulness of its own audit trail.

This paper presents an agentic AI framework for intelligent patient monitoring and clinical
decision support designed around these three shortfalls. The framework organizes specialized
LLM agents --- monitoring, diagnosis, risk prediction, treatment recommendation, explanation, and
verification --- under a coordinator, over a memory layer that treats the patient's longitudinal
EHR timeline as a first-class retrieval corpus, with a dedicated verification gate and an audit
trail whose faithfulness is itself an evaluation target. Clinician oversight is structured
rather than advisory: recommendations reach the clinician only after verification, with linked
evidence and calibrated confidence, a design choice supported by meta-analytic evidence that
unstructured human--LLM collaboration yields fragile gains \cite{wang2026collaboration}.

The contributions are: (1) a layered framework that grounds retrieval in the patient timeline
rather than only in external knowledge; (2) a verification gate and evidence-linked audit trail
specified as measurable components, with metrics for recommendation grounding and trail
faithfulness; (3) an evaluation design on MIMIC-IV that scores longitudinal patient tracking
rather than one-shot question answering, positioned against the exam-based and synthetic-EHR
benchmarks that dominate current practice \cite{jiang2025medagentbench,arora2025healthbench,lovon2025mimic}; and (4) an analysis of how the framework's safeguards map onto emerging
regulatory requirements for unconfined non-deterministic clinical software \cite{tan2026undcs};
and (5) a pilot feasibility study on the openly licensed MIMIC-IV demo that grounds the
framework's distinguishing mechanisms in real ICU data --- timelines for 140 stays build in
under three seconds, timestamp-aware retrieval answers in roughly 12 ms without returning
future information, the verification gate's evidence requirement suppresses false alerts
faster than true ones (86\% versus 57\% blocked at four required signals), and every logged
audit-trail reference re-resolves against the source record. These are feasibility results
on a 100-patient cohort, and the paper reports them with that caveat throughout.

The remainder of the paper reviews related work (Section~II), details the framework and its
evaluation design (Section~III), discusses implications and limitations (Section~IV), and
reports the pilot study (Section~V).


\section{Related Work}

\subsection{LLM Agents and Their Infrastructure}

The agent paradigm grew from reasoning-and-acting loops \cite{yao2023react} and self-taught tool use
\cite{schick2023toolformer} into general architectures with memory, planning, and reflection
\cite{wang2024survey,shinn2023reflexion}. Multi-agent frameworks --- AutoGen's conversational
orchestration \cite{wu2024autogen}, CAMEL's role-play \cite{li2023camel}, MetaGPT's SOP-driven pipelines
\cite{hong2024metagpt} --- established coordination patterns that recent surveys have formalized into
taxonomies of collaboration types and structures \cite{tran2025collaboration} and, at the systems
level, into brain-inspired modular blueprints \cite{liu2025foundationagents}. Agent memory has its
own literature organizing construction, management, and retrieval operations \cite{wu2025memory}, and
interoperability is standardizing around MCP for tool access and A2A-style protocols for
inter-agent exchange \cite{ehtesham2025protocols}. Sapkota et al. distinguish single tool-using
agents from orchestrated agentic ecosystems and identify healthcare as a primary beneficiary of
the latter \cite{sapkota2025agents}. This infrastructure maturity is what makes our framework
buildable; none of it, however, addresses what a clinical agent should be grounded in or how its
outputs should be verified.

\subsection{Medical LLMs and Medical Agents}

Medical LLMs progressed from encoding clinical knowledge \cite{singhal2023clinical} through
expert-level question answering \cite{singhal2025medpalm2} to multimodal generalists \cite{tu2024generalist}
and open alternatives \cite{toma2023clinicalcamel}. The agentic turn followed: MedAgents assembles
expert-role panels for zero-shot reasoning \cite{tang2024medagents}; Agent Hospital evolves doctor
agents in a simulated hospital \cite{li2024agenthospital}; AgentClinic embeds them in simulated
dialogue \cite{schmidgall2024agentclinic}; EHRAgent executes code over structured records
\cite{shi2024ehragent}. The 2025 generation is more capable: TxAgent reasons over 211 therapeutic
tools \cite{gao2025txagent}, RiskAgent routes to validated risk calculators \cite{liu2025riskagent},
MedRAX orchestrates chest X-ray tools in a training-free loop \cite{fallahpour2025medrax},
DoctorAgent-RL learns proactive consultation policies \cite{feng2025doctoragent}, and the AMIE line
extends conversational diagnosis to longitudinal management and multimodal evidence
\cite{palepu2025disease,saab2025multimodal}. Surveying this landscape, Wang et al. conclude that
clinical decision-making attracts the most research and the least deployment evidence
\cite{wang2025baymax}. Across all of these systems, patient state is reconstructed from the prompt or
scenario at each step; none maintains a persistent, evidence-linked model of one real patient
across an admission.

\subsection{Retrieval-Augmented Generation in Medicine}

RAG is the standard mitigation for hallucination \cite{lewis2020rag,gao2023rag}, refined by
self-critical variants \cite{asai2024selfrag} and, recently, folded into the agent loop as agentic
RAG --- retrieval that is planned, executed iteratively, and reflected upon \cite{singh2025agenticrag}.
In medicine, MedRAG couples retrieval with a diagnostic knowledge graph \cite{zhao2025medrag}, and a
ten-model npj Digital Medicine study shows guideline-grounded RAG can exceed human accuracy in a
bounded preoperative task \cite{ke2025ragfitness}. The corpus, in every case, is external knowledge.
Retrieval anchored to the patient's own longitudinal record --- the grounding a monitoring system
actually needs --- remains unexplored, which is the specific opening this work occupies.

\subsection{Evaluation, Safety, and Regulation}

Evaluation practice is catching up with agent capabilities. MedAgentBench places agents in a
FHIR-compliant virtual EHR with physician-authored tasks \cite{jiang2025medagentbench}; HealthBench
grades open-ended health conversations against physician rubrics \cite{arora2025healthbench}; general
surveys catalog agent-evaluation methodology and flag safety as under-measured
\cite{yehudai2025evaluation}. On real records, revisited MIMIC-IV baselines show text-based models
competitive with tabular classifiers for one-shot prediction \cite{lovon2025mimic}, and COMPOSER-LLM
demonstrates prospective sepsis alerting with an LLM adjudicating high-uncertainty cases
\cite{shashikumar2025sepsis}. The safety literature supplies taxonomies of medical hallucination
\cite{kim2025hallucinations}, arguments for calibrated uncertainty \cite{atf2025uncertainty}, threat
models for trustworthy agents \cite{yu2025trustagent}, and adversarial stress tests showing that
multi-agent topology shapes harm propagation \cite{chen2025medsentry}. Regulatory analyses find that
unconstrained LLMs already produce medical-device-like output \cite{weissman2025unregulated} and
propose governance for unconfined non-deterministic clinical software \cite{tan2026undcs}, while the
only meta-analysis of clinician--LLM collaboration reports fragile gains and clinically
significant error rates \cite{wang2026collaboration}. What this literature demands --- grounded
recommendations, calibrated confidence, inspectable trails, structured oversight --- is precisely
what no current system measures end-to-end; our framework is designed so that these properties
are evaluated rather than claimed.


\section{Methodology}

\subsection{Design Overview}

The framework follows a design-science approach \cite{hevner2004design}: an artifact (the layered
agentic framework) is constructed to close identified gaps and evaluated against explicit
metrics. Six layers organize the system --- data, memory, reasoning and knowledge, agent
orchestration, clinical decision, and trustworthy AI --- with a clinician dashboard and
human-in-the-loop validation closing the loop. The full architecture is specified in the thesis
(Chapter 3); this section states the components that carry the paper's claims.

\subsection{Patient-Timeline Retrieval}

The memory layer treats the patient's longitudinal EHR trajectory as a first-class retrieval
corpus, alongside --- not instead of --- external medical knowledge. Four stores are maintained:
short-term working context; long-term patient memory (admissions, diagnoses, medications,
procedures); a vector database over clinical notes and events; and clinical context memory for
the active monitoring episode. Memory operations follow the construction / management /
retrieval decomposition of the agent-memory literature \cite{wu2025memory}: writes are explicit
events linked to source records, management handles decay and revision (a corrected lab value
supersedes but does not erase its predecessor), and retrieval is timestamp-aware so that
evidence reflects what was knowable at decision time. Retrieval-augmented generation then draws
on both corpora --- the patient timeline and guideline/literature knowledge
\cite{lewis2020rag,singh2025agenticrag} --- so that every recommendation can cite patient-specific
evidence, not only population-level knowledge \cite{zhao2025medrag,ke2025ragfitness}.

\subsection{Agent Orchestration}

A coordinator agent delegates to specialized agents --- monitoring, planner, diagnosis, risk
prediction, treatment recommendation, explanation, and verification --- using ReAct-style
reasoning within each agent \cite{yao2023react}. Coordination follows a hub-and-spoke topology: all
inter-agent messages pass through the coordinator, a choice motivated by evidence that
decentralized medical multi-agent topologies propagate adversarial or erroneous influence
further \cite{chen2025medsentry}. Tool access uses MCP-style standardized interfaces and inter-agent
exchange follows A2A-style conventions \cite{ehtesham2025protocols}, keeping the prototype
composable with hospital tooling. External tools follow the validated-tool principle of
RiskAgent and TxAgent --- calculations and lookups are delegated to auditable components rather
than generated \cite{liu2025riskagent,gao2025txagent}.

\subsection{Verification Gate and Audit Trail}

The verification agent is the framework's distinguishing safety component. Every candidate
recommendation is checked for (a) evidential grounding --- is the claim entailed by the retrieved
patient evidence and guidelines, screened against known classes of medical hallucination
\cite{kim2025hallucinations}; (b) guideline compliance and conflict with active orders; and (c) a
calibrated confidence estimate that accompanies the recommendation to the clinician
\cite{atf2025uncertainty,guo2017calibration}. Recommendations failing any check are blocked and
returned with the failure reason. Every check writes to an evidence-linked audit trail; the
trail's faithfulness --- whether it accurately reflects the evidence the system actually used ---
is treated as a measured property, not an assumed one, using entailment-based spot audits
\cite{es2024ragas}. This combination is designed to satisfy, and to demonstrate compliance with,
emerging expectations for unconfined non-deterministic clinical software: guardrails,
moderation, retrieval grounding, and inspectability \cite{tan2026undcs,weissman2025unregulated}.

\subsection{Human-in-the-Loop Protocol}

Clinicians receive only verified recommendations, each with linked evidence, confidence, and
explanation. Review is structured --- approve, modify, or reject, with reasons captured into
memory --- rather than free-form, reflecting meta-analytic evidence that unstructured
clinician--LLM collaboration produces fragile gains \cite{wang2026collaboration}. Rejections and
modifications feed the monitoring loop as supervision signals.

\subsection{Evaluation Design}

Evaluation uses MIMIC-IV \cite{johnson2023mimic} cohorts (sepsis and deterioration use cases;
Sepsis-3 labels \cite{singer2016sepsis3,vincent1996sofa}) and scores the system on longitudinal
tracking, not one-shot prediction. Primary axes: (1) decision quality against clinical ground
truth and revisited MIMIC-IV baselines \cite{lovon2025mimic}; (2) grounding rate --- the fraction of
recommendation claims entailed by retrieved patient evidence, following RAG-evaluation practice
\cite{es2024ragas}; (3) verification-gate effectiveness --- ungrounded-recommendation catch rate at
fixed review budget; (4) audit-trail faithfulness under spot audit; and (5) rubric-graded
recommendation quality in the style of physician-rubric benchmarks
\cite{arora2025healthbench,jiang2025medagentbench}. Ablations remove patient-timeline retrieval,
the verification gate, and longitudinal memory in turn, isolating each claimed contribution.
Statistical comparisons follow standard practice for correlated classifiers
\cite{delong1988comparing,dietterich1998approximate} with false-discovery control
\cite{benjamini1995controlling}.


\section{Discussion}

\subsection{What the Framework Claims and What It Does Not}

The framework's novelty is deliberately narrow. Tool-using medical agents exist
\cite{gao2025txagent,liu2025riskagent,fallahpour2025medrax}; multi-agent coordination exists
\cite{wu2024autogen,tran2025collaboration}; clinical RAG exists and works within its grounding
assumptions \cite{zhao2025medrag,ke2025ragfitness}. We do not claim these components. We claim the
three things the 2025--2026 literature documents as absent: retrieval grounded in the patient's
own timeline, evaluation that scores longitudinal tracking on real ICU records, and a
verification gate whose audit trail is itself measured. If the ablations show these components
do not improve grounded decision quality, the framework's thesis fails honestly --- the
components are separable and testable by design.

\subsection{Relation to Contemporary Evaluation Practice}

MedAgentBench and HealthBench moved evaluation toward realism --- synthetic FHIR records and
physician rubrics respectively \cite{jiang2025medagentbench,arora2025healthbench} --- and the best
agents complete only about seventy percent of virtual-EHR tasks, which counsels humility about
bedside readiness. Our evaluation design borrows their instruments (task framing, rubric
grading) but changes the substrate to real, longitudinal MIMIC-IV admissions
\cite{johnson2023mimic,lovon2025mimic}. The trade-off is acknowledged: MIMIC-IV is retrospective
and single-region, and retrospective evaluation cannot establish the prospective safety that
COMPOSER-LLM's deployment begins to demonstrate for one task \cite{shashikumar2025sepsis}.
Retrospective longitudinal evaluation on real records is nonetheless the missing rung between
exam benchmarks and prospective trials, and it is the rung this work supplies.

\subsection{Safety, Oversight, and Regulation}

The safety design responds to documented failure modes rather than hypothetical ones: medical
hallucinations reach practice \cite{kim2025hallucinations}; uncommunicated uncertainty erodes
appropriate reliance \cite{atf2025uncertainty}; multi-agent topologies can amplify a single
compromised agent \cite{chen2025medsentry}; and unconstrained LLM output already meets medical-device
criteria \cite{weissman2025unregulated}. Structured human oversight is a load-bearing component, not
a disclaimer --- the meta-analytic evidence shows collaboration gains are fragile when oversight
is unstructured \cite{wang2026collaboration}. Mapping the framework's guardrails onto proposed UNDCS
requirements \cite{tan2026undcs} is an early attempt to make a research prototype legible to the
regulatory conversation; we expect the mapping, not the architecture, to need revision as rules
mature.

\subsection{Limitations}

Four limitations bound the claims. First, empirical evidence so far is a pilot on the
100-patient MIMIC-IV demo without the LLM loop (Section~V); the full retrospective evaluation
is designed but not executed, and no claim of prospective clinical benefit is made. Second,
MIMIC-IV's critical-care population limits generalization to other care settings
\cite{johnson2023mimic}. Third, verification quality is
bounded by the entailment methods used to measure it \cite{es2024ragas}; a verification gate can
only be as trustworthy as its own evaluation, which is why trail faithfulness is spot-audited
by humans as well. Fourth, LLM API costs and latency constrain the monitoring cadence the
prototype can sustain, a practical ceiling the agent-evaluation literature identifies as
generally under-reported \cite{yehudai2025evaluation}.

\subsection{Future Work}

Beyond executing the Chapter 4 evaluation, three directions follow. Prospective shadow-mode
deployment in the style of COMPOSER-LLM is the natural next evidence level
\cite{shashikumar2025sepsis}. Federated or on-premises open models \cite{toma2023clinicalcamel} would
address the privacy constraints that cloud LLMs impose on real EHR timelines. And the
verification gate invites formalization: moving from entailment heuristics toward auditable,
possibly symbolic, checking of clinical claims against structured evidence --- the direction the
trustworthy-agent literature identifies as open \cite{yu2025trustagent,tan2026undcs}.


\section{Pilot Feasibility Study}

\subsection{Rationale and Scope}

The full evaluation of Section~III-F requires credentialed MIMIC-IV access and an LLM agent
loop; both are in progress. To establish that the framework's distinguishing components are
implementable and behave as designed on real ICU data, we conducted a bounded pilot on the
openly licensed MIMIC-IV Clinical Database Demo (v2.2; 100 patients, 275 admissions, 140 ICU
stays) \cite{johnson2023mimic}. The pilot exercises the non-LLM slice of the architecture ---
timeline construction, timestamp-aware retrieval, an alert-generating classical baseline, and
the verification gate with its audit trail --- and is reported as feasibility evidence. With
twenty deaths in the cohort, no performance claim survives this sample size, and we make none.

\subsection{Setup}

For each ICU stay, the memory layer assembled a chronological timeline from vitals, laboratory
results, medication starts, and strictly prior admissions, each event carrying a provenance
reference to its source table and row. Retrieval at a decision point (24 h after ICU
admission) scored events by recency, type, and deterioration relevance, and never returned an
event recorded after the query time. A standardized logistic regression over first-24 h
features (5-fold stratified cross-validation) produced an in-hospital mortality risk score,
and the top quintile generated 28 alerts (7 in patients who died, 21 in patients who
survived). The verification gate passed an alert only when the patient's retrieved timeline
contained at least \emph{m} distinct deterioration-relevant signals (tachycardia, hypotension,
tachypnea, hypoxemia, temperature derangement, elevated lactate, leukocytosis or leukopenia,
rising creatinine or bilirubin, thrombocytopenia) within the six hours preceding the decision.
Every gate decision logged its evidence references, and trail faithfulness was measured by
re-resolving each logged reference against the raw tables.

\subsection{Results}

Timeline construction processed all 140 stays in 2.9 s (median 438 events per stay,
IQR 236--798); 43\% of stays carried at least one prior admission, confirming that longitudinal
context exists to retrieve even in a small cohort. Retrieval answered queries in 11.7 ms
median (17.9 ms p95) with full top-k coverage and returned no future information by
construction. The baseline reached a cross-validated AUROC of 0.641 (95\% bootstrap CI
0.478--0.792) --- an interval that crosses chance, as expected at this sample size; its role is
to supply a realistic alert stream, not predictive skill.

The gate produced the pilot's most instructive result. As the required number of distinct
recent signals rises from 1 to 6, the pass rate for false alerts falls from 0.90 to 0.05
while the pass rate for true alerts falls more slowly (0.57 to 0.14); at m = 4 the gate
blocks 86\% of false alerts while retaining 43\% of true ones. Evidence-gating therefore
preferentially suppresses unsupported alerts, which is the mechanism the framework relies on.
The same experiment exposes the mechanism's cost: three of the seven true alerts had no
abnormal signal in the six-hour window at all, so the gate imposes a sensitivity floor
governed by the evidence window --- a safety parameter to be tuned, and an effect invisible in
exam-style evaluation. All 183 logged evidence references re-resolved exactly against the raw
tables (trail resolvability 1.00); in this deterministic pilot the perfect score is expected,
and the metric exists precisely because it becomes non-trivial once an LLM writes the trail.

\subsection{What the Pilot Does and Does Not Show}

The pilot shows that the framework's memory, retrieval, and verification mechanisms are
implementable on real ICU records at interactive latency, that evidence-gating discriminates
in the intended direction, and that audit-trail faithfulness is a measurable quantity. It
does not show clinical utility, generalization beyond a 100-patient demo cohort, or anything
about the LLM agent loop, which the full MIMIC-IV evaluation of Section~III-F will test. The
demo lacks the clinical notes module, so retrieval here is structured-data only.


\section{Conclusion}

The 2025--2026 literature settled the question of whether LLM agents can act in clinical
settings; it left open whether their actions can be trusted for a specific patient over
time. This paper's answer is architectural and methodological rather than model-centric: a
layered multi-agent framework whose distinguishing commitments --- the patient timeline as a
retrieval corpus, a verification gate that checks recommendations against retrieved patient
evidence, an audit trail whose faithfulness is measured, and structured rather than advisory
clinician oversight --- target exactly the properties the field's own benchmarks and
regulatory analyses identify as absent. A pilot on the open MIMIC-IV demo demonstrates that
the non-LLM substrate of these commitments is implementable at interactive latency and that
evidence-gating discriminates in the intended direction, blocking unsupported alerts at a
substantially higher rate than supported ones while exposing an honest sensitivity floor that
exam-style evaluation cannot reveal. The full evaluation on credentialed MIMIC-IV, with the
LLM agent loop in place, is designed and pending; it, not the pilot, will test the clinical
substance of the claims. If the framework's components fail to improve grounded decision
quality there, the design permits that failure to be located precisely --- which is, in a
domain where unverifiable success is indistinguishable from plausible error, the property we
argue clinical agentic AI needs most.
